%

%

%%%%

%%%   Самарский государственный технический университет

%%%   Кафедра прикладной математики и информатики

%%%

%%%   Преамбула для статьи в журнал "Вестник Самарского государственного

%%%   технического университета. Серия: «Физико-математические науки»"

%%%   Ориентирована под MikTEx 2.4 for Win32

%%%

%%%   Хотя сам журнал собирается под GNU/Linux на дистрибутиве TeXLive



\documentclass[11pt,a4paper,twoside]{article}



\usepackage[cp1251]{inputenc}

\usepackage[T2A]{fontenc}

\usepackage[english,russian]{babel}

\usepackage{samgtubib}

\usepackage[dvips]{graphicx}

\usepackage[multiple]{footmisc}



\usepackage{samgtu}

\renewcommand{\VestnikYear}{2009} % Год издания Вестника

\renewcommand{\VestnikNumber}{2\,(19)} % Номер Вестника

\renewcommand{\VestnikMonth}{Ноябрь} % Месяц выхода Вестника

\renewcommand{\VestnikData}{01.11.09} % Дата подписания Вестника в печать

\renewcommand{\VestnikUPL}{26,64} % Усл. печ. л.

\renewcommand{\VestnikUIL}{25,73} % Уч.-изд. л.

\renewcommand{\VestnikRegNumber}{479/09} % Регистрационный номер

\renewcommand{\VestnikOrder}{994} % Номер заказа







%\usepackage{pst-barcode}

%\usepackage{auto-pst-pdf}

%

%\newcommand{\totop}[1]{\raisebox{6pt}[1pt][2pt]{#1}} 

%\newcommand{\tottop}[1]{\raisebox{12pt}[1pt][2pt]{#1}} 

%\newcommand{\totttop}[1]{\raisebox{18pt}[1pt][2pt]{#1}} 

%\newcommand{\tottttop}[1]{\raisebox{24pt}[1pt][2pt]{#1}} 

%\newcommand{\totttttop}[1]{\raisebox{30pt}[1pt][2pt]{#1}} 

%\newcommand{\tobot}[1]{\raisebox{-6pt}[1pt][2pt]{#1}} 

%\newcommand{\tobbot}[1]{\raisebox{-12pt}[1pt][2pt]{#1}} 

%\newcommand{\tobbbot}[1]{\raisebox{-18pt}[1pt][2pt]{#1}}

%

%%\samgtupubl % для генерации pdf, отдаваемого в типографию СамГТУ

%\pdfcrop % обрезка pdf для размещения на math-net.ru

%\draft % На корректуре

%%\nofilllastpage % Последняя русскоязычная страница заполнена не полностью

%\filllastpage

%%\briefcomm % Статья является кратким сообщением



\vsgtu{1332}



\FirstPage{31}

\LastPage{43}









%\begin{document}





\renewcommand{\baselinestretch}{0.94}





\hfill \begin{pspicture}(1in,1in)

\psbarcode{http://dx.doi.org/10.14498/vsgtu1332}{}{qrcode}

\end{pspicture}

\vspace{-1in}



\udk{517.956.27}

\msc{35R30, 35J25}



\rutitle{Об определении неизвестных коэффициентов при~старших производных в линейном эллиптическом уравнении}% Полный заголовок

{Об определении неизвестных коэффициентов \dots}



\ruauthor{Р.~А.~Алиев}{А\,л\,и\,е\,в~Р.~А.}



\ruaffil{\ruazuk.}

\enaffil{\enazuk.}







\ruauthordetails{1}{\fullauthor{\rupid{52318}{Рамиз Аташ оглы Алиев}} \regalia{к.ф.-м.н.,

доц.; \mem{ramizaliyev3@rambler.ru}}\position{доцент} \dept{каф. информатики}}



\enauthordetails{1}{\fullauthor{\enpid{52318}{Ramiz A. Aliyev}} \regalia{Cand. Phys. \& Math. Sci.; \mem{ramizaliyev3@rambler.ru}} \position{Associate Professor} \dept{Dept. of Informatics}}



\rukeywords{обратная задача, эллиптическое уравнение, метод последовательных приближений, существование и единственность решения}



\ruabstract{Исследуется обратная задача нахождения коэффициентов и решения линейного эллиптического уравнения в заданном прямоугольнике. Доказана теорема существования, единственности и устойчивости решения поставленной обратной задачи. С помощью метода последовательных приближений построен регуляризирующий алгоритм для определения нескольких коэффициентов.}



\entitle{On the determination of the unknown coefficients of the highest derivatives in a linear elliptic equation}



\enauthor{R.~A.~Aliyev}{A\,l\,i\,y\,e\,v~R.~A.}











\enabstract{Inverse problems on restoration of coefficients to the differential equations with partial derivatives are of interest in many applied researches. These problems lead to necessity of the approached decision of inverse problems for the equations of mathematical physics which are incorrect in classical sense. In the article the existence, uniqueness and stability of the solution of the given inversion problem for the elliptic equation are proved.}



\enkeywords{inverse problem, elliptic equation, successive approximations method, existence and uniqueness of solutions}



\date{16/VI/2014}

\revisiondate{11/VII/2014}

\accepteddate{27/VIII/2014}



\makerutitle







\Section[n]{Введение} 

Обратные задачи по определению коэффициентов дифференциальных уравнений с частными производными представляют интерес во многих прикладных исследованиях \cite{ali:bib:1, ali:bib:2}. 

Эти задачи приводят к необходимости приближенного решения обратных задач математической физики, которые некорректны в классическом смысле. К ним относятся задачи идентификации неизвестных плотностей источников и коэффициентов уравнения. 

Большое значение имеют коэффициентные задачи для эллиптических уравнений, в~которых неизвестные коэффициенты не зависят от одной переменной. 

Такие модели характерны для задач теории фильтрации. 

В частности, определение теплофизических характеристик сред в стационарном случае приводит к обратным задачам для эллиптических уравнений [\citen{ali:bib:4}--\citen{ali:bib:10}]. 

Отметим также цикл работ [\citen{ali:bib:11}--\citen{ali:bib:23}], посвященный теоретическому и численному исследованию обратных задач для линейных и нелинейных уравнений эллиптического типа. 

Исследование таких задач вызвано как теоретическим интересом, так и практической необходимостью.



Целью данной работы является доказательство единственности и существования решений обратной краевой задачи для эллиптического уравнения второго порядка с дополнительным условием.



\smallskip



\Section{Постановка задачи}

Рассмотрим задачу определения функций $\{a_1(x_2)$, $a_2(x_2)$, $u(x_1,x_2)\}$ из следующих условий: 

\begin{gather}

\label{ali:1}

-a_1(x_2)u_{x_1x_1}-a_2(x_2)u_{x_2x_2}+c(x_2)u=h(x_1,x_2), \quad (x_1,x_2)\in D;

\\

\label{ali:2}

u(0,x_2)=\phi_1(x_2), \quad u(l_1,x_2)=\phi_2(x_2), \quad 0\le x_2\le l_2;

\\

\label{ali:3}

u(x_1,0)=\varphi_1(x_1), \quad u(x_1,l_2)=\varphi_2(x_1), \quad 0\le x_1\le l_1;

\\

\label{ali:4}

a_1(x_2)u_{x_1}(0,x_2)=g_1(x_2), \quad 0\le x_2\le l_2;

\\

\label{ali:5}

a_2(x_2)u_{x_1}(l_1,x_2)=g_2(x_2), \quad 0\le x_1\le l_2,

\end{gather}

причём $\varphi_1(0)=\phi_1(0)$, $\varphi_2(l_1)=\phi_2(l_2)$, $\varphi_1(l_1)=\phi_2(0)$, $\varphi_2(0)=\phi_1(l_2)$. 

Здесь $D=\{(x_1,x_2)\mid 0<x_1<l_1, 0<x_2<l_2\}$; 

$c(x_2)\in C^\alpha[0, l_2]$, 

$h(x_1,x_2)$, 

$h_{x_1 x_1}(x_1,x_2)\in C^\alpha(\overline D)$, 

$g_i(x_2)\in C^\alpha[0, l_2]$, 

$\phi_i(x_2)\in C^{2+\alpha}(0, l_2)\cap C[0, l_2]$, 

$\varphi_i(x_1)\in C^{2+\alpha}(0, l_1)\cap C[0, l_1]$ "--- заданные функции; $i=1, 2$; $0<\alpha<1$.



\smallskip

\begin{definition}[1]Функции $\{a_1(x_2)$, $a_2(x_2)$, $u(x_1,x_2)\}$ назовём решением задачи \eqref{ali:1}--\eqref{ali:5}, если $u(x_1,x_2)\in C^2(D)\cap C(\bar{D})$, $0<a_i(x_2)\in C[0, l_2]$, $i=1, 2$,  и удовлетворяются соотношения \eqref{ali:1}--\eqref{ali:5}.

\end{definition}



\smallskip



Нетрудно проверить, что если решение задачи \eqref{ali:1}--\eqref{ali:5} существует, то при принятых предположениях о гладкости данных в задаче $a_i(x_2)\in C^\alpha[0, l_2]$, $i=1, 2$, $u(x_1, x_2)\in C^{2+\alpha}(\overline D)$. 

Действительно, при принятых предположениях из общей теории эллиптических уравнений следует, 

что $u(x_1,x_2)\in W_p^2(D)\subset C^{1+\alpha}(\bar{D})$ при $p>2$ \cite[c.~283]{ali:bib:24}. 

Поэтому из дополнительных условий \eqref{ali:4} и \eqref{ali:5} следует, что $a_i(x_2)\in C^\alpha[0, l_2]$, $i=1, 2$. 

Поэтому $u(x_1,x_2)\in C^{2+\alpha}(\overline D)$. 



\smallskip



\Section{Единственность и устойчивость решения}

Пусть кроме задачи \eqref{ali:1}--\eqref{ali:5} задана еще задача $(\bar{\text{\ref{ali:1}}})$--$(\bar{\text{\ref{ali:5}}})$, где все функции, входящие в \eqref{ali:1}--\eqref{ali:5}, заменены соответствующими функциями с чертой. 

Положим 

$$Z(x_1, x_2)=\bar u(x_1,x_2)-u(x_1,x_2), \quad \lambda_i(x_2)=\bar{a}_i(x_2)-a_i(x_2),$$ 

$$\delta_1(x_2)=\bar{c}(x_2)-c(x_2),\quad  \delta_2(x_1,x_2)=\bar h(x_1,x_2)-h(x_1,x_2),$$ 

$$\delta_{i+2}(x_2)=\bar{\phi}_i(x_2)-\phi_i(x_2), \quad \delta_{i+4}(x_1)=\bar{\varphi}_i(x_1)-\varphi_i(x_1),$$ 

$$\delta_{i+6}(x_2)=\bar{g}_i(x_2)-g_i(x_2 ),\quad i=1, 2.$$



Через $\tilde{\delta}_2(x_1,x_2)$ обозначим функцию на границе, совпадающую соответственно с $\delta_{i+2}(x_2)$, $\delta_{i+4}(x_1)$, $i=1, 2$, и принадлежащую $C^{2+\alpha}(\bar{D})$.



\smallskip



\hypertarget{ali:lemma1}{}



\begin{lemma}[1]Пусть решения задачи \eqref{ali:1}--\eqref{ali:5} существуют{\rm .} 

Тогда верны следующие 

оценки{\rm:}

$$|u(x_1,x_2)|\le\max\Bigl[\max\limits_D\Bigl|\frac{h(x_1,x_2)}{c(x_2)}\Bigr|,\max\limits_i \max\limits_{x_1}|\varphi_i(x_1)|,\max\limits_i \max\limits_{x_2}  |\phi_i(x_2)|\Bigr],$$

\begin{multline*}

|u_{x_1 x_1}(x_1,x_2)|\le\max\Bigl[\max\limits_D \Bigl|\frac{h_{x_1 x_1}(x_1,x_2)}{c(x_2)}\Bigr|, \max\limits_i \max\limits_{x_1}|\varphi_{i_{x_1 x_1}}(x_1)|, \\ \max\limits_{x_2}|\theta_{10}(0,x_2)|,\max\limits_{x_2}|\theta_{20}(l_1,x_2)|\Bigr],

\end{multline*}

\begin{multline*}

|u_{x_2 x_2}(x_1,x_2)|\le\max\Bigl[\max\limits_D \frac{1}{a_2(x_2)}|a_1(x_2)u_{x_1 x_1}+c_1(x_2)u(x_1,x_2)-h(x_1,x_2)|,\\

\max\limits_i \max\limits_{x_1} |\phi_{ix_1 x_1}|,  \max_{x_1}|\theta_{11}(x_1,0)|,\max\limits_{x_1}|\theta_{21}(x_1,l_2)|\Bigr];

\end{multline*}







\noindent

здесь

\begin{multline*}

\theta_{i k}(x_1,x_2)=\frac{1}{a_{k+1}(x_2)}\Bigl\{(1-k)

\bigl[-a_2(x_2)\phi_{ix_2x_2}(x_2)+c(x_2)\phi_i(x_2) \bigr]+ \\ +

 k \bigl[-a_1(x_2)\varphi_{ix_1x_1}(x_1)+c(x_2)\varphi_i(x_1)\bigr]- h(x_1,x_2)\Bigr\},\quad  i = 1, 2, \; k = 0, 1.

\end{multline*}

\end{lemma}



\begin{proof} Первое неравенство получается из принципа максимума. Уравнение \eqref{ali:1} продифференцируем дважды по $x_1$, учитывая условия \eqref{ali:3} и обозначая $\upsilon(x_1, x_2)=u_{x_1x_1}$, получим 

\begin{gather*}

-a_1(x_2)\upsilon_{x_1x_1}-a_2(x_2)\upsilon_{x_2x_2}+c(x_2)u=h_{x_1x_1}(x_1,x_2), \quad (x_1,x_2)\in D;

\\

\upsilon(0,x_2)=\theta_{10}(0,x_2), \quad \upsilon(l_1,x_2)=\theta_{20}(l_1,x_2), \quad 0\le x_2\le l_2;

\\

\upsilon(x_1,0)=\varphi_{1_{x_1x_1}}(x_1), \quad u(x_1,l_2)=\varphi_{2_{x_1x_1}}(x_1), \quad 0\le x_1\le l_1.

\end{gather*}



Используя принцип максимума, получим вторую оценку. Последнее неравенство получается из уравнения. 

Лемма \hyperlink{ali:lemma1}{1} доказана. \end{proof}



\smallskip



Единственность решения обратной задачи \eqref{ali:1}--\eqref{ali:5} в предположении его существования устанавливает следующая теорема.



\smallskip



\hypertarget{ali:thm1}{}



\begin{theorem}[1]Пусть $g_1(x_2)\ne0,$ $g_2(x_2)\ne0,$ $N\mathop{\rm mes}D<1.$ 

Тогда решение задачи \eqref{ali:1}--\eqref{ali:5} единственно и верна следующая оценка{\rm :}

\begin{multline}

\sum\limits_{i=1}^2 

\bigl\|\bar{a}_i(x_2)-a_i(x_2)\bigr\|_{C[0,l_2]}+

\bigl\|\bar{u}-u\bigr\|_{C(\overline D)}\le 

N_1\biggl[ \bigl\|\bar{c}(x_2)-c(x_2)\bigr\|_{C[0,l_2]}+

\\

+\bigl\|\bar{h}(x_1,x_2)-h(x_1,x_2)\bigr\|_{C(\overline D)}+\sum\limits_{i=1}^2\bigl(

\bigl\|\bar{\phi}_i(x_2)-\phi_i(x_2)\bigr\|_{C^2[0,l_2]}+

\\

+\bigl\|\bar{\varphi}_i(x_1)-\varphi_i(x_1)\bigr\|_{C^2[0,l_1]}\bigr)+\sum\limits_{i=1}^2 \bigl \|\bar{g}_i(x_2)-g_i(x_2)\bigr\|_{C[0,l_2]}\biggr],

\label{ali:7}

\end{multline}

где  $N,$ $N_1$ "--- положительные постоянные{\rm ,} зависящие от данных и решения задачи.

\end{theorem}



\smallskip

\begin{proof} Из $(\bar{\text{\ref{ali:1}}})$--$(\bar{\text{\ref{ali:5}}})$ соответственно вычтем \eqref{ali:1}--\eqref{ali:5} и положим $Z_1(x_1,x_2)=Z(x_1,x_2)-\tilde{\delta}_2(x_1,x_2)$. 

Тогда получим 

\begin{multline}

\label{ali:8}

-\bar{a}_1(x_2)Z_{1_{x_1x_1}}-\bar{a}_2(x_2)Z_{1_{x_2x_2}}= \\

=\delta_9(x_1,x_2)-\bar{c}(x_2)\tilde{\delta}_2(x_1,x_2)-\bar{c}(x_2)Z_1 +\sum\limits_{i=1}^2{\alpha_i}(x_1,x_2)\lambda_i(x_2);

\end{multline}

\begin{gather}

\label{ali:10}

Z_1(0,x_2)=0, \quad Z_1(l_1,x_2)=0;

\\

\label{ali:10}

Z_1(x_1,0)=0, \quad Z_1(x_1,l_2)=0;

\\

\label{ali:11}

\lambda_i(x_2)=\delta_{i+9}(x_2)+\gamma_i(x_2) Z_{1x_1}[(i-1)l_1, x_2], \quad i=1, 2. 

\end{gather}

Здесь 

$$\alpha_i(x_1,x_2)=u_{x_ix_i},\quad  \delta_9(x_1,x_2)=-\delta_1(x_2)u+\delta_2(x_1,x_2)+\sum\limits_{i=1}^2\bar{a}_i(x_2)\tilde{\delta}_{2_{x_ix_i}}(x_1,x_2),$$ $$\gamma_i(x_2)=\bar{a}_i(x_2)\{-u_{x_1}[(i - 1)l_1,x_2]\}^{-1},$$

$$\delta_{i+9}(x_2)= \gamma_i(x_2) \bigl\{ [-\bar{a}_i(x_2)]^{-1}\delta_{i+6}(x_2)    +

\tilde \delta _{2 x_1} [(i-1)l_1, x_2 ]

  \bigr\} , \quad i=1, 2.$$



Обозначим

\begin{gather*}

L_{ik}(x_j)=\Bigl(\frac{x_j}{i}\Bigr)^{k(j-1)}\frac{(2-i)l_j-(-1)^{i+1}({1}/{2})^{k(j-1)}x_j}{l_j},

\\

P_i(x_j)=(-1)^{(j+1)i}l_j^{(3-2i)(j-1)}x_j^{(i-1)(j-1)}, \quad X_j=x_{j+(-1)^{j+1}},

\\

t_{ijk}=(1-j+k)\delta_{6-k}[(i-1)l_1]+[(-1)^k(1-j)+1-k]\delta_{(i+2)x_2}(kl_2), 

\end{gather*}

$i,j = 1, 2$, $k = 0, 1$.



Функция $\tilde{\delta}_2(x_1,x_2)$ определяется следующим образом:

\begin{multline*}

\tilde{\delta}_2(x_1,x_2)=\sum\limits_{i=1}^2\biggl\{\sum\limits_{j=1}^2L_{i0}(x_j)\bigl\{\delta_{i+2j}(X_j)+(1-j)L_{i0}(x_1)\delta_{i+2j}(0)\bigr\}+

\\

+(-1)^{i+1}\frac{x_1x_2}{l_1l_2}\delta_4[(i-1)l_2]\biggr\}-\frac{x_1}{l_1}\delta_4(0).

\end{multline*}

При помощи функции Грина \cite{ali:bib:25} из \eqref{ali:8}--\eqref{ali:10} определим функцию $Z_1(x_1,x_2)$ через правую часть равенства и это выражение подставим в условие \eqref{ali:11}. 

Тогда получим

\begin{equation}

\begin{array}{l}

\displaystyle

Z(x_1,x_2)=\tilde{\delta}_2(x_1,x_2)+\int_D G(x_1,x_2,\xi_1,\xi_2)\biggl[\delta_9(\xi_1,\xi_2)-\bar{c}(\xi_2)Z(\xi_1,\xi_2)+

\\

\displaystyle \hspace{6cm}

+\sum\limits_{i=1}^2 \alpha_i(\xi_1,\xi_2)\lambda_i(\xi_2)\biggr]d\xi_1d\xi_2,

\\

\displaystyle

\lambda_i(x_2)=\delta_{i+9}(x_2)+\gamma_i(x_2)\int_D G_{x_1}[(i-1)l_1,x_2,\xi_1,\xi_2]\biggl[\delta_9(\xi_1,\xi_2)-

\\

\displaystyle \hspace{2.5cm}

-\bar{c}(\xi_2)Z(\xi_1,\xi_2)+\sum\limits_{i=1}^2{\alpha_i}(\xi_1,\xi_2)\lambda_i(\xi_2)\biggr]d\xi_1d\xi_2, \quad i = 1, 2.

\end{array}

\label{ali:12}

\end{equation}



Функция Грина имеет следующие оценки \cite{ali:bib:25}:

\begin{gather}

|G(x_1,x_2,\xi_1,\xi_2)|\le M_1\ln\frac{1}{\sqrt{(x_1-\xi_1)^2+(x_2-\xi_2)^2}},

\\

\label{ali:13}

|G_{x_1}(x_1,x_2,\xi_1,\xi_2)|\le M_2[(x_1-\xi_1)^2+(x_2-\xi_1)^2]^{- 1/2}, \quad M_i>0,\; i=1, 2.

\end{gather}



При достаточно малой мере $D$ для указанных интегралов в \eqref{ali:12} су\-щест\-вуют 

оценки $N_i[\mathop{\rm mes}D]^{1/2}$, $N_i>0$, $i=2, 3, 4$. 

Теперь в системе \eqref{ali:12} положим 

$$

\chi=\max\limits_{x_1,x_2}|Z(x_1,x_2)|+\sum\limits_{i=1}^2 \max\limits_{x_2}|\lambda_i(x_2)|.

$$



Из системы \eqref{ali:12} получим

\begin{multline*}

\chi\le N_5\biggl[\bigl\|\delta_1(x_2)\bigr\|_{C[0,l_2]}+

\bigl\|\delta_1(x_1,x_2)\bigr\|_{C(\overline{D})}+

\bigl\|\tilde{\delta}_2(x_1,x_2)\bigr\|_{C^2(\bar{D})}+

\\ 

+\sum\limits_{i=1}^3 \bigl\|\delta_{i+6}(x_2)\bigr\|_{C[0,l_2]}\biggr]+\chi N_6[\mathop{\rm mes}D]^{1/2},

\end{multline*}

где

$N_5$, $N_6$ "--- некоторые положительные числа. 

Обозначим $N=N_6^2$. По условию \hyperlink{ali:thm1}{теоремы}  $N\mathop{\rm mes}D<1$, поэтому $N_6[mesD]^{1/2}<1$. 

Отсюда получим, что при $(x_1,x_2)\in\bar{D}$ верна оценка устойчивости \eqref{ali:7}. Единственность решения задачи следует из оценки \eqref{ali:7}. Теорема  \hyperlink{ali:thm1}{1} доказана. \end{proof}



\smallskip



\Section{Метод последовательных приближений}

Метод последовательных приближений для решения задачи \eqref{ali:1}--\eqref{ali:5} применяется по следующей схеме:

\begin{gather}

\label{ali:14}

-a_1^{(s)}(x_2)u_{x_1x_1}^{(s+1)}-a_2^{(s)}(x_2)u_{x_2x_2}^{(s+1)}+c(x_2)u^{(s+1)}{=}h(x_1,x_2), \; (x_1,x_2)\in D;

\\

\label{ali:16}

u^{(s+1)}(0,x_2)=\phi_1(x_2), \quad u^{(s+1)}(l_1,x_2)=\phi_2(x_2), \quad 0\le x_2\le l_2;

\\

\label{ali:16}

u^{(s+1)}(x_1,0)=\varphi_1(x_1), \quad u^{(s+1)}(x_1,l_2)=\varphi_2(x_1), \quad 0\le x_1\le l_1;

\\

\label{ali:17}

a_1^{(s+1)}(x_2)u_{x_1}^{(s+1)}(0,x_2)=g_1(x_2), \quad 0\le x_2\le l_2;

\\

\label{ali:18}

a_2^{(s+1)}(x_2)u_{x_1}^{(s+1)}(l_1,x_2)=g_2(x_2), \quad 0\le x_2\le l_2.

\end{gather}



По схеме \eqref{ali:14}--\eqref{ali:18} последовательные итерации проводятся следующим образом. 

Сперва выбираются некоторые $a_{i}^{(0)}(x_2)>0$, $i=1, 2$, принадлежащие $C^\alpha[0,l_2]$, и подставляются в уравнение \eqref{ali:14}. 

Далее решается задача \eqref{ali:14}--\eqref{ali:16} и находится $u^{1}(x_1,x_2)$. 

По функциям $u_{x_1}^{(1)}(0,x_2)$, $u_{x_1}^{(1)}(l_1,x_2)$ из условий \eqref{ali:17}--\eqref{ali:18} находятся $a_1^{(1)}(x_2)$, $a_2^{(1)}(x_2)$ и эти функции используются для проведения следующего шага итерации. 



\smallskip



\hypertarget{ali:thm2}{}



\begin{theorem}[2]Пусть решение задачи \eqref{ali:1}--\eqref{ali:5} существует и при 

всех $s=0, 1$  и т$.$д$.$ $u^{(s)}(x_1,x_2)\in C^2(D),$ 

$a_{i}^{(s)}(x_2)\in C^\alpha[0, l_2],$ $i=1, 2,$ 

${g_1(x_2)u_{x_1}^{(s)}(0,x_2)>0,}$ 

$g_2(x_2)u_{x_1}^{(s)}(l_1,x_2)>0,$ $N\mathop{\rm mes}D<1,$ а

производные $u^{(s)}(x_1,x_2)$ по $x_1,$ $x_2$ до второго порядка равномерно ограничены. 

Тогда функции $\{a_1^{(s)}(x_2),$ $a_2^{(s)}(x_2),$ $u^{(s)}(x_1,x_2)\},$ полученные методом последовательных приближений \eqref{ali:14}--\eqref{ali:18}{\rm ,} при $s\to+\infty$ равномерно сходятся к~решению задачи \eqref{ali:1}--\eqref{ali:5} со скоростью геометрической прогрессии$.$ $N$ "--- положительная постоянная{\rm ,} зависящая от данных задачи{\rm .}

\end{theorem}



\smallskip



\begin{proof} Положим 

$$

Z^{(s)}(x_1,x_2)=u(x_1,x_2)-u^{(s)}(x_1,x_2), \; \lambda_{i}^{(s)}(x_2)=a_i(x_2)-a_i^{(s)}(x_2), \; i=1, 2.

$$

Легко проверить, что эти функции удовлетворяют системе 

\begin{gather}

\label{ali:19}

-a_1(x_2)Z_{x_1x_1}^{(s+1)}-a_2(x_2)Z_{x_2x_2}^{(s+1)}+c(x_2)Z^{(s+1)}{=}\sum\limits_{i=1}^2\alpha_i^{(s)}(x_1,x_2)\lambda_i^{(s)}(x_2);

\\

\label{ali:21}

Z^{(s+1)}(0,x_2)=0, \quad Z^{(s+1)}(l_1,x_2)=0,

\\

\label{ali:21}

Z^{(s+1)}(x_1,0)=0, \quad Z^{(s+1)}(x_1,l_2)=0,

\\

\label{ali:22}

\lambda_i^{(s+1)}(x_2)=\gamma_i^{(s)}(x_2)Z_{x_1}^{(s+1)}[(i-1)l_1,x_2], \quad i=1, 2,

\end{gather}

где 

$

\alpha_{i_0}^{(s)}(x_1,x_2)=u_{x_ix_i}^{(s+1)}, \; \gamma_i^{(s)}(x_2)=a_i(x_2)\{-u_{x_1}^{(s+1)}[(i-1)l_1,x_2]\}^{-1}, \; i=1, 2.

$



С помощью функции Грина из \eqref{ali:19}--\eqref{ali:21} определим $Z^{(s+1)}(x_1,x_2)$ через правую часть равенства \eqref{ali:19} и подставим это выражение в условия \eqref{ali:22}. 

Тогда получим

\begin{multline}

\label{ali:23}

\lambda_i^{(s+1)}(x_2)=\gamma_i^{(s)}(x_2)\int_D G_{x_1}\bigl[(i-1)l_1,x_2,\xi_1,\xi_2\bigr]\times \\

\times \sum\limits_{i=1}^2\alpha_i^{(s)}(\xi_1,\xi_2)\lambda_i^{(s)}(\xi_2)d\xi_1d\xi_2,  \quad 

i = 1, 2.

\end{multline}



Положим 

$$

\chi^{(s)}=\sum\limits_{i=1}^2\max\limits_{x_2}\bigl|\lambda_i^{(s)}(x_2)\bigr|

.$$

 Аналогично из системы \eqref{ali:23} следует, что  $\chi^{(s+1)}\le\chi^{(s)}N_6[\mathop{\rm mes}D]^{1/2}$. Таким образом, теорема \hyperlink{ali:thm2}{2} доказана.~\end{proof}

 

 \smallskip



\Section{Существование решения}

Пусть $m(x_2)\in C^2[0,l_2]$, $m(x_2)>0$, $m^{\prime}(0)>0$, $m^{\prime}(l_2)<0$, $m^{\prime\prime}(x_2)\ge0$. 

Введём следующие обозначения: 

\begin{gather*}

\eta(x_1)=\frac{3}{2}\mu_0^{-1}(x_1-l_1)^2,\quad  \beta_{i+2n}=m^{(n)}(il_2-l_2),

\\

\xi_{in}(x_1)=\eta(x_1)\varphi_i(l_1)+\beta_{i+2n}l_1^{-1}(l_1-x_1),\\

p_i(x_i)=(2-i)\xi_{10}(x_1)+(i-1)\xi_{21}(x_1),\\

M_1=\max\limits_{x_2}l_1^{-1}[\phi_1(x_2)-\phi_2(x_2)],\quad  

M_{i+1}=\max[M_1,\max\limits_{x_1}|\varphi_{ix_1}(x_1)|],\\

M=\max\limits_{k=2, 3}M_k,\quad  i, j=1, 2, \; n=0, 1.

\end{gather*}



\smallskip



\hypertarget{ali:lemma2}{}



\begin{lemma}[2]Пусть задача

\begin{gather*}

-a_1(x_2)u_{x_1x_1}-a_2(x_2)u_{x_2x_2}+u=0, \quad (x_1,x_2)\in D;

\\

u(0,x_2)=\phi_1(x_2), \quad u(l_1,x_2)=\phi_2(x_2), \quad 0\le x_2\le l_2;

\\

u(x_1,0)=\varphi_1(x_1), \quad u(x_1,l_2)=\varphi_2(x_1), \quad 0\le x_1\le l_1,

\end{gather*}

удовлетворяющая условиям 

$$\varphi_1(0)=\phi_1(0),\quad \varphi_1(l_1)=\phi_2(0),\quad \varphi_2(l_1)=\phi_2(l_2),\quad \phi_1(l_2)=\varphi_2(0),$$ 

при заданных $a_1(x_2),$ $a_2(x_2)\ge\mu_0>0$ имеет решение{\rm ,} которое  принадлежит классу $C^2(D)\cap C(\overline D),$ и выполняются условия 

$$0<l_1<1,\quad \mu_0^{-1}l_1<1,\quad m(x_2)\le\frac{1}{2}\phi_2(x_2)\le\frac{1}{3}\bigl[M-l_1^{-1}m(x_2)\bigr],$$ 

$$\beta_i\le\frac{1}{2}\varphi_i(l_1)\le\frac{1}{3}\bigl(M-\beta_i l_1^{-1}\bigr),$$ $$m(x_2)+\eta(0)\phi_2(x_2)\le\phi_1(x_2)-\phi_2(x_2)\le Ml_1,$$ 

$$x_1 l_1^{-1}\beta_i\le\varphi_i(0)-\varphi_i(x_1)\le Mx_1,$$ 

$$\xi_{i0}(x_1)\le\varphi_i(x_1)-\varphi_i(l_1)\le M(l_1-x_1),$$ 

$$\phi_{ix_2x_2}(x_2)=0,\quad i=1, 2,$$ 

тогда

\begin{equation}

\label{ali:24}

\begin{array}{c}

-M-\phi_i(x_2)(2\mu_0)^{-1}l_1\le u_{x_1}(il_1-l_1,x_2)\le-m(x_2)l_1^{-1}, \\ i=1, 2, \quad  0<x_2<l_2.

\end{array}

\end{equation}

\end{lemma}





\begin{proof} Положим

$$

\upsilon_i(x_1,x_2)=u(x_1,x_2)+m(x_2)l_1^{-1}(x_1-il_1+l_{1})-\phi_i(x_2)-(i-1)\eta(x_1)\phi_i(x_2),

$$

\begin{multline*}

V_i(x_1,x_2)=-u(x_1,x_2)+\phi_i(x_2)-M(x_1-il_1+l_1)+

\\

+(-1)^i(2\mu_0)^{-1}x_1(l_1-x_1)\phi_i(x_2), \quad i = 1,2.

\end{multline*}



Нетрудно проверить, что $\upsilon_1(x_1,x_2)$ удовлетворяет условиям задачи

\begin{gather*}

-a_1(x_2)\upsilon_{1{x_1x_1}}-a_2(x_2)\upsilon_{1{x_2x_2}}+\upsilon_1= \hspace{4cm} \\

\hspace{4cm} = -\phi_1(x_2)+m(x_2)x_1l_1^{-1}- a_2 (x_2) m^{\prime\prime}(x_2)x_1l_1^{-1},

\\

\upsilon_1(0, x_2)=0, \quad \upsilon_1(l_1,x_2)=-\phi_1(x_2)+m(x_2)+\phi_2(x_2),

\\

\upsilon_1(x_1, 0)=\varphi_1(x_1)-\varphi_1(0)+m(0)x_1l_1^{-1}, \\ \upsilon_1(x_1,l_2)=\varphi_2(x_1)-\varphi_2(0)+m(l_2)x_1l_1^{-1}.

\end{gather*}



Поэтому, по условию \hyperlink{ali:lemma2}{леммы}, наибольшее положительное значение функции $\upsilon_1(x_1,x_2)$ достигается при $x_1=0$. 

Тогда $\upsilon_{1_{x_1}}(0,x_2)\le0$, другими словами,

\begin{equation}

\label{ali:25}

u_{x_1}(0,x_2)\le-m(x_2)l_1^{-1}.

\end{equation}



Аналогично прежнему, подставляя $V_1(x_1,x_2)$ и учитывая условия леммы, получаем, что наибольшее положительное значение функции $V_1(x_1,x_2)$ достигается при $x_1=0$. 

Поэтому $V_{1_{x_1}}(0,x_2)\le0$ или

\begin{equation}

\label{ali:26}

-M-\phi_1(x_2)(2\mu_0)^{-1}l_1\le u_{x_1}(0, x_2).

\end{equation}



Объединяя оценки \eqref{ali:25} и \eqref{ali:26}, получим оценку \eqref{ali:24} при $i=1$. 

Аналогично прежнему получим оценку \eqref{ali:24} при $i=2$. Лемма \hyperlink{ali:lemma2}{2} доказана. 

\end{proof}



\smallskip



\hypertarget{ali:thm3}{}



\begin{theorem}[3]Пусть 

$$\phi_i(x_2)\in C^{2+\alpha}(0, l_2)\cap C[0, l_2],\quad \varphi_i(x_1)\in C^{2+\alpha}(0, l_1)\cap C[0, l_1],$$ 

$$0<l_1<1, \quad \mu_0^{-1}l_1<1,\quad c(x_2)=1,\quad h(x_1,x_2)=0,$$  

$$m(x_2)\le\frac{1}{2}\phi_2(x_2)\le\frac{1}{3}\bigl[M- l_1^{-1} m(x_2)\bigr],\quad 

\beta_i\le\frac{1}{2}\varphi_i(l_1)\le\frac{1}{3}\bigl(M-\beta_il_1^{-1}\bigr),$$ $$m(x_2)+\eta(0)\phi_2(x_2)\le\phi_1(x_2)-\phi_2(x_2)\le Ml_1,$$ 

$$x_1l_1^{-1}\beta_i\le\varphi_i(0)-\varphi_i(x_1)\le Mx_1,\quad \xi_{i0}(x_1)\le\varphi_i(x_1)-\varphi_i(l_1)\le M(l_1-x_1),$$ 

$$g_i(x_2)<0,\quad g_0\le-g_i(x_2)-\frac{1}{2}\phi_i(x_2)l_1,$$ 

$$\varphi_{ix_1}(jl_1-l_1)<0,\quad \phi_{ix_2x_2}(x_2)=0,\quad i, j=1, 2;$$ 

$m(x_2)$ "--- неотрицательная функция такая$,$ что $g_i(x_2)[m(x_2)]^{-1},$ $i=1, 2,$ огра\-ни\-ченны$;$ $g_0$ "--- положительное число{\rm .} Тогда задача \eqref{ali:1}--\eqref{ali:5} имеет хотя бы одно решение{\rm .}

\end{theorem} 



\begin{proof} Нетрудно видеть, что

$$

a_j(il_2-l_2)=g_j(il_2-l_2)[\varphi_{ix_1}(jl_1-l_1)]^{-1}, \quad i,j = 1, 2.

$$



Доказательство проводится методом последовательных приближений. 

Из утверждения леммы \hyperlink{ali:lemma2}{2} следует, что 

$$

-M\bigl[1+\phi_1(x_2)(2g_0)^{-1}l_1\bigr]\le u_{x_1}^{(s+1)}(il_1-l_1,x_2)\le-m(x_2)\cdot l_1^{-1},

$$

$i = 1, 2,$  $0<x_2<l_2,$

тогда 

$$

g_0M^{-1}\le a_i^{(s+1)}(x_2)\le\max\limits_{x_2}\bigl\{[-g_i(x_2)]\cdot[m(x_2)]^{-1}\bigr\}\cdot l_1, \; i=1, 2, \; 0<x_2<l_2.

$$



Таким образом, при всех приближениях $a_1^{(s)}(x_2^{(s)})$, $a_2^{(s)}(x_2^{(s)})$ "--- строго положительные, 

непрерывные и равномерно ограниченные функции. 

Тогда из общей теории эллиптических уравнений следует, что при условиях теоремы последовательность $\{u^{(s)}(x_1,x_2)\}$ равномерно ограничена по норме $W_p^2(D)$ $\forall p>2$. 

Поэтому $\{u^{(s)}(x_1,x_2)\}$ компактна в $C^1(\overline D)$. При этом из условий \eqref{ali:17}, \eqref{ali:18} следует, что $\{a_1^{(s)}(x_2),$ $a_2^{(s)}(x_2)\}$ будет компактна в $C[0,l_2]$. 

Отсюда и из \eqref{ali:14}--\eqref{ali:16} вытекает компактность $\{u^{(s)}(x_1,x_2)\}$ в $C^2(\overline D)$. 

В~системе \eqref{ali:4}--\eqref{ali:18}, переходя к пределу при $s\to+\infty$, получим, что существует пара функции $\{a_1(x_2),a_2(x_2),u(x_1,x_2)\}$, удовлетворяющая условиям \eqref{ali:1}--\eqref{ali:5}.

 Теорема \hyperlink{ali:thm3}{3} доказана. \end{proof}

 

 \smallskip





Вместо условия \eqref{ali:3} можно выбрать одно из следующих условий:

\begin{gather}

\label{ali:27}

u_{x_2}(x_1,0)=\varphi_1(x_1), \quad u_{x_2}(x_1,l_2)=\varphi_2(x_1), \quad 0\le x_1\le l_1;

\\

\label{ali:28}

u_{x_2}(x_1,0)=\varphi_1(x_1), \quad u(x_1,l_2)=\varphi_2(x_1), \quad 0\le x_1\le l_1;

\\

\label{ali:29}

u(x_1,0)=\varphi_1(x_1), \quad u_{x_2}(x_1,l_2)=\varphi_2(x_1), \quad 0\le x_2\le l_1. 

\end{gather}

Для таких задач $\tilde{\delta}_2(x_1,x_2) $ определяются соответственно в следующем виде:

$$

\begin{array}{l}

\displaystyle

\tilde{\delta}_2(x_1,x_2)=\sum\limits_{i,j=1}^2L_{i1}(x_j)\delta_{i+2j}(X_j)+ \\ 

\displaystyle

\hspace{2cm}+ (-1)^{j+1}x_2\Bigl(\frac{x_2}{2l_2}+j-2\Bigr)L_{i1}(x_1)\delta_{{(i+2)}_{x_2}}\bigl[(j - 1)l_2\bigr],

\end{array}

$$

$$

\begin{array}{l}

\displaystyle

\tilde{\delta}_2(x_1,x_2)=\sum\limits_{i,j=1}^2P_i(x_j)L_{i0}(x_j)\delta_{i+2j}(X_j)+l_2^{2-j}t_{ij0}L_{i0}(x_1)[L_{j0}(x_2)]^{2^{j-1}},

\\

\displaystyle

\tilde{\delta}_2(x_1,x_2)=\sum\limits_{i,j=1}^2l_j^{(i-1)(j-1)}\bigl[L_{i0}(x_j)\bigr]^{2^{(2-i)(j-1)}}\delta_{i+2j}(X_j)- \\ 

\displaystyle

\hspace{7cm}

-l_2^{j-1}t_{ij1}L_{i0}(x_1)\bigl[L_{j0}(x_2)\bigr].

\end{array}

$$





Аналогично теореме \hyperlink{ali:thm3}{3} можно доказать следующие теоремы существования для соответствующих задач. 







\begin{theorem}[4]Пусть 

$$\phi_i(x_2)\in C^{1+\alpha}(0,l_2)\cap C[0, l_2], \quad \varphi_i(x_1)\in C^{2+\alpha}(0,l_1)\cap C[0,l_1],$$

$$0<l_1<1,\quad \mu_0^{-1}l_1<1, \quad c(x_2)=1, \quad h(x_1,x_2)=0,\quad \varphi_1(0)\le0,\quad \varphi_2(0)\ge0,$$ $$m(x_2)\le\frac{1}{2}\phi_2(x_2)\le\frac{1}{3}[M- l_1^{-1} m(x_2)],$$ 

$$-\frac{2}{3}\beta_{i+2}l_1^{-1}(2-i)\le\varphi_i(l_1)\le-\frac{2}{3}\beta_{i+2}l_1^{-1}(i-1),$$ $$\eta(0)\phi_2(x_2)+m(x_2)\le\phi_1(x_2)-\phi_2(x_2)\le M_1l_1,$$ 

$$x_1l_1^{-1}\beta_{i+2}(2i-3)\le\varphi_i(2x_1-ix_1)-\varphi_i(ix_1-x_1)\le 0,$$ 

$$(i-1)\xi_{i1}(x_1)\le\varphi_i(x_1)-\varphi_i(l_1)\le(2-i)\xi_{i1}(x_1),$$ 

$$g_i(x_2)<0, \quad g_0\le-g_i(x_2)-\frac{1}{2}\phi_i(x_2)l_1,\quad  \phi_{ix_2x_2}(x_2)=0,\quad i=1, 2;$$ 

$m(x_2)$ "--- неотрицательная функция такая{\rm ,} что $g_i(x_2)[m(x_2)]^{-1},$ $i=1, 2,$ ограниченны{\rm ;} 

$g_0$ "--- положительное число{\rm .} 

Тогда задача {\rm \eqref{ali:1}--\eqref{ali:2}{\rm ,} \eqref{ali:4}--\eqref{ali:5} } и \eqref{ali:27} имеет хотя бы одно решение{\rm .} 

\end{theorem}



\smallskip



\begin{theorem}[5]Пусть 

$$\phi_i(x_2)\in C^{2+\alpha}(0,l_2)\cap C[0,l_2], \quad \varphi_1(x_1)\in C^{1+\alpha}(0,l_1)\cap C[0,l_1],$$ 

$$\varphi_2(x_1)\in C^{2+\alpha}(0,l_1)\cap C[0, l_1], \quad 0<l_1<1,\quad \mu_0^{-1}l_1<1,$$ 

$$c(x_2)=1, \quad h(x_1,x_2)=0, \quad c(x_2)=1, \quad \varphi_1(0)\le0, \quad \varphi_2(0)\ge0,$$ 

$$m(x_2)\le\frac{1}{2}\phi_2(x_2)\le\frac{1}{3}[M_3- l_1^{-1} m(x_2)],$$ 

$$-\frac{2}{3}\beta_{i+2}l_1^{-1}(2-i)\le\varphi_i(l_1)\le-\frac{2}{3}\beta_{i} \bigl[ M_3 - l_1^{-1} \beta_i \bigr] (i-1),$$ 

$$\eta(0)\phi_2(x_2)+m(x_2)\le\phi_1(x_2)-\phi_2(x_2)\le M_3l_1,$$ 

$$(-1)^i\beta_{({i+2})/{i}}x_1l_1^{-1}\le\varphi_i(2x_1-ix_1)-\varphi_i(ix_1-x_1)\le M_3x_1(i-1),$$ 

$$\xi_{i0}(x_1)(i-1)\le\varphi_i(x_1)-\varphi_i(l_1)\le\xi_{i1}(x_1)(2-i)+M_3(l_1-x_1)(i-1),$$ 

$$g_i(x_2)<0,\quad g_0\le-g_i(x_2)-\frac{1}{2}\phi_i(x_2)l_1, \quad \varphi_{2x_1}(il_1-l_1)<0,$$ 

$$\phi_{ix_2x_2}(x_2)=0, \quad i=1, 2;$$  

$m(x_2)$ "--- неотрицательная функция такая{\rm ,} что $g_i(x_2)[m(x_2)]^{-1},$ $i=1, 2,$ огра\-ни\-ченны$;$ $g_0$ "--- положительное число{\rm .} Тогда задача {\rm \eqref{ali:1}--\eqref{ali:2}{\rm ,} \eqref{ali:4}--\eqref{ali:5}} и \eqref{ali:28} имеет хотя бы одно решение{\rm .}

\end{theorem} 







\begin{theorem}[6]Пусть 

$$\phi_i(x_2)\in C^{2+\alpha}(0, l_2)\cap C[0, l_2], \quad \varphi_1(x_1)\in C^{2+\alpha}(0, l_1)\cap C[0,l_1],$$ $$\varphi_2(x_1)\in C^{1+\alpha}(0, l_1)\cap C[0, l_1],\quad 0<l_1<1, \quad \mu_0^{-1}l_1<1,$$ 

$$c(x_2)=1, \quad h(x_1,x_2)=0,\quad 0\le\varphi_i(l_1)\le\frac{2}{3}[-\beta_{i^2}l_1^{-1}+M_2(2-i)],$$

$$p_i(x_1)\le\varphi_i(x_1)-\varphi_i(l_1)\le M_2(l_1-x_1)(2-i),$$ 

$$x_1l_1^{-1}\beta_{i^2}\le\varphi_i(0)-\varphi_i(x_1)\le(2-i)M_2x_1,$$ 

$$\eta(0)\phi_2(x_2)+m(x_1)\le\phi_1(x_2)-\phi_2(x_2)\le M_2l_1 ,$$ 

$$g_i(x_1)<0,\quad g_0\le-g_i(x_2)-\frac{1}{2}\phi_i(x_2)l_1,\quad \varphi_{1x_1}(il_1-l_1)<0,$$ 

$$\phi_{ix_2x_2}(x_2)=0,\quad i=1, 2;$$ 

$m(x_2)$ "---  неотрицательная функция такая$,$ что $g_i(x_1)[m(x_2)]^{-1},$ $i=1, 2,$ огра\-ниченны$;$ $g_0$ "--- положительные числа{\rm .} Тогда задача {\rm \eqref{ali:1}--\eqref{ali:2}{\rm ,} \eqref{ali:4}--\eqref{ali:5}} и \eqref{ali:29} имеет хотя бы одно решение{\rm .}

\end{theorem} 





%\thanks{{\bf ORCID}

%

%{\bf Ramiz Aliyev:} \url{http://orcid.org/xxxx-xxxx-xxxx-xxxx}

%}







\RusRef



\begin{thebibliography}{99}







\RBibitem{ali:bib:1} 

\by Лаврентьев~М.~М., Романов~В.~Г., Шишатский~С.~П.

\book Некорректные задачи математической физики и анализа

\publaddr М.

\publ Наука 

\yr 1980

\totalpages 288



\RBibitem{ali:bib:2} 

\by Кабанихин~С.~И.

\book Обратные и некорректные задачи

\publaddr М.

\publ Наука 

\yr 2009

\totalpages 458





%\RBibitem{ali:bib:3} 

%\by Искендеров~А.~Д. 

%\paper Об одной обратной задаче для уравнения эллиптического типа

%\jour Всесоюз. симпозиум по обратным задачам для дифференциальных уравнений

%\yr 1972

%\pages 107--115





\RBibitem{ali:bib:4} 

\by Искендеров~А.~Д. 

\paper Некоторые обратные задачи об определении правых частей дифференциальных уравнений 

\jour Изв. Акад. наук Азерб. ССР

\yr 1976 

\issue 2 

\pages 58--63

\zmath{https://zbmath.org/?q=an:0345.35084}

%Iskenderov, A.D.

%Some inverse problems on the determination of right-hand parts of differential equations. (Russian) Zbl 0345.35084

%Izv. Akad. Nauk Az. SSR, Ser. Fiz.-Tekh. Mat. Nauk 1976, No.2, 58-63 (1976).





\RBibitem{ali:bib:5} 

\by Искендеров~А.~Д.

\paper Обратная задача об определении коэффициентов квазилинейного эллиптического уравнения 

\jour Изв. Акад. наук Азерб. ССР

\yr 1978 

\issue 2 

\pages 80--85

\zmath{https://zbmath.org/?q=an:0411.35088}

%Iskenderov, A.D.

%Inverse problem about the determination of the coefficients of a quasilinear elliptic equation. (Russian) Zbl 0411.35088

%Izv. Akad. Nauk Az. SSR, Ser. Fiz.-Tekh. Mat. Nauk 1978, No.2, 80-85 (1978).





\RBibitem{ali:bib:6} 

\by  Искендеров~А.~Д. 

\paper Обратная задача об определении коэффициентов эллиптического уравнения 

\jour Диффер. уравн.

\yr 1979

\vol 15

\pages 858--867

\zmath{https://zbmath.org/?q=an:0404.35094}

%Iskenderov, A.D.

%The inverse problem of the determination of the coefficients in an elliptic equation. (English) Zbl 0428.35083

%Differ. Equations 15, 605-612 (1979).





\RBibitem{ali:bib:7} 

\by Темирбулатов~С.~И.

\book Обратные задачи для эллиптических уравнений 

\publaddr Алма-Ата

\publ Казах. ун-т

\yr 1975 

\totalpages 72





\RBibitem{ali:bib:8} 

\by Клибанов~М.~В.

\paper Двумерная обратная задача для одного эллиптического уравнения

\jour Диффер. уравн.

\yr 1983

\vol 19

\issue 6 

\pages 1072--1074

\zmath{https://zbmath.org/?q=an:0529.35075}

%Klibanov, M.V.

%A two-dimensional inverse problem for an elliptic equation. (Russian) Zbl 0529.35075

%Differ. Uravn. 19, No.6, 1072-1074 (1983).







\RBibitem{ali:bib:9} 

\by Клибанов~М.~В.

\paper Обратные задачи в целом и карлемановские оценки 

\jour Диффер. уравн.

\yr 1984

\vol 20

\issue 6 

\pages 1035--1041

\zmath{https://zbmath.org/?q=an:0573.35083}

%Klibanov, M.V.

%Inverse problems in the ”large” and Carleman bounds. (English. Russian original) Zbl 0573.35083

%Differ. Equations 20, 755-760 (1984); translation from Differ. Uravn. 20, No.6, 1035-1041 (1984).





\RBibitem{ali:bib:10} 

\by Клибанов~М.~В.

\paper Единственность в целом обратных задач для одного класса дифференциального уравнения

\jour Диффер. уравн.

\yr 1984

\vol 20

\issue 11

\pages 1947--1953

\zmath{https://zbmath.org/?q=an:0567.35080}

%Klibanov, M.V.

%Uniqueness in the large of solutions of inverse problems for a class of differential equations. (English. Russian original) Zbl 0567.35080

%Differ. Equations 20, 1390-1395 (1984); translation from Differ. Uravn. 20, No.11, 1947-1953 (1984).





\RBibitem{ali:bib:11} 

\by Хайдаров~А.

\paper Один класс обратных задач для эллиптических уравнений 

\jour Докл. Акад. наук СССР

\yr 1984

\vol 277

\issue 6 

\pages 1335--1337

% One class of inverse problems for elliptic equations 

%Authors of Document Khaidarov, A.

%Year the Document was Publish 1984

%Source of the Document Dokl. Akad. Nauk SSSR

%277, pp. 1335-1337







\RBibitem{ali:bib:12} 

\by Хайдаров~А.

\paper Об одной обратной задаче для эллиптических уравнений

\inbook Некоторые задачи математической физики и анализа

\publ Наука. Сибирское отделение

\publaddr Новосибирск

\yr 1984

\pages 245--249

%A. Khaidarov, “On an inverse problem for elliptic equations,” in: Some Problems of Mathematical Physics and Analysis [in Russian], Nauk, Novosibirsk (1984), pp. 245–249.





\RBibitem{ali:bib:13} 

\by Вабищевич~П.~Н.

\paper Обратная задача восстановления правой части эллиптического уравнения и ее численные решения

\jour Диффер. уравн.

\yr 1985

\vol 21

\issue 2

\pages 277--284

\zmath{https://zbmath.org/?q=an:0576.65122}

%Vabishchevich, P.N.

%An inverse problem for reconstructing the right-hand side of an elliptic equation and its numerical solution. %\zmath{https://zbmath.org/?q=an:0595.65130}

%Differ. Equations 21, 201-207 (1985).



\RBibitem{ali:bib:14} 

\by Вабищевич~П.~Н.

\paper О единственности некоторых обратных задач для эллиптических уравнений

\jour Диффер. уравн.

\yr 1988

\vol 24

\issue 12

\pages 2125--2129

\zmath{https://zbmath.org/?q=an:0695.35214}

%Vabishchevich, P.N.

%Uniqueness of solutions of some inverse problems for elliptic equations. (English. Russian original) Zbl 0695.35214

%Differ. Equations 24, No.12, 1443-1446 (1988); translation from Differ. Uravn. 24, No.12, 2125-2129 (1988).





\RBibitem{ali:bib:15} 

\by Соловьев~В.~В.

\zmath{https://zbmath.org/?q=an:1160.35549}

\paper Обратные задачи для эллиптических уравнений на плоскости.~I

\jour Диффер. уравн.

\yr 2006

\vol 42

\issue 8

\pages 1106--1114

%\zmath{https://zbmath.org/?q=an:1160.35549}

%Inverse problems for elliptic equations on the plane. I. (English. Russian original) Zbl 1160.35549

%Differ. Equ. 42, No. 8, 1170-1179 (2006); translation from Differ. Uravn. 42, No. 8, 1106-1114 (2006).

%\crossref{10.1134/s0012266106080118}

%\scopus{2-s2.0-33750032371}



\RBibitem{ali:bib:16} 

\by Соловьев~В.~В.

\paper Обратные задачи определения источника и коэффициента в~эллиптическом уравнении в~прямоугольнике

\jour Ж. вычисл. матем. и матем. физ.

\yr 2007

\vol 47

\issue 8

\pages 1365--1377

\mathnet{http://mi.mathnet.ru/zvmmf265}

\mathscinet{http://www.ams.org/mathscinet-getitem?mr=2378182}

%\transl

%Solov??v, V. V. 

%Source and coefficient inverse problems for an elliptic equation in a rectangle

%\jour Comput. Math. Math. Phys.

%\yr 2007

%\vol 47

%\issue 8

%\pages 1310--1322

%\mathscinet{http://www.ams.org/mathscinet-getitem?mr=2378182}

%\crossref{http://dx.doi.org/10.1134/S096554250708009X}

%\scopus{2-s2.0-34548411765}





\Bibitem{ali:bib:17} 

\paper A Global Uniqueness Theorem for an Inverse Boundary Value Problem

\jour Annals of Mathematics 

\vol 125 

\issue 1 

\pages 153--169

\by Sylvester~J., Uhlmann~G.

\crossref{10.2307/1971291}



\Bibitem{ali:bib:18} 

\paper Inverse coefficient problems for nonlinear elliptic equations

\jour ANZIAM Journal 

\vol 49 

\yr 2008

\issue 02 

\pages  271--279

\by Yang~R., Ou~Y.

\crossref{10.1017/s1446181100012839}

\scopus{2-s2.0-77954661742}





\RBibitem{ali:bib:19} 

\by Вахитов~И.~С.

\paper Обратная задача идентификации старшего коэффициента в уравнении диффузии-реакции

\jour Дальневост. матем. журн.

\yr 2010

\vol 10

\issue 2

\pages 93--105

\mathnet{http://mi.mathnet.ru/dvmg62}

\elib{http://elibrary.ru/item.asp?id=15484861}





\RBibitem{ali:bib:20} 

\by Козлов~В.~А., Мазья~В.~Г., Фомин~А.~В.

\paper Об одном итерационном методе решения задачи Коши для эллиптических уравнений

\jour Ж. вычисл. матем. и матем. физ.

\yr 1991

\vol 31

\issue 1

\pages 64--74

\mathnet{http://mi.mathnet.ru/zvmmf3145}

\mathscinet{http://www.ams.org/mathscinet-getitem?mr=1099360}

\zmath{https://zbmath.org/?q=an:0733.65056}

%\transl

%Kozlov, V. A.; Maz?ya, V. G.; Fomin, A. V. An iterative method for solving the Cauchy problem for elliptic equations

%\jour U.S.S.R. Comput. Math. Math. Phys.

%\yr 1991

%\vol 31

%\issue 1

%\pages 45--52

%\mathscinet{http://www.ams.org/mathscinet-getitem?mr=1099360}

%\isi{http://gateway.isiknowledge.com/gateway/Gateway.cgi?GWVersion=2&SrcApp=PARTNER_APP&SrcAuth=LinksAMR&DestLinkType=FullRecord&DestApp=ALL_WOS&KeyUT=A1991HU21600006}





\Bibitem{ali:bib:21} 

\paper An iterative procedure for solving a Cauchy problem for second order elliptic equations

\yr 2004 

\jour Mathematische Nachrichten 

\vol 272 

\issue 1 

\pages 46--54

\by Johansson~T.

\crossref{10.1002/mana.200310188} 

\scopus{2-s2.0-4644351704}







\RBibitem{ali:bib:22} 

\by Алиев~Р.~А.

\paper Теорема единственности одной обратной задачи для квазилинейного уравнения эллиптического типа

\jour Известия вузов. Северо-Кавказский регион. Естественные науки

\yr 2012

\issue 5 

\pages 5--10

\elib{http://elibrary.ru/item.asp?id=18051841}



\RBibitem{ali:bib:23}

\by Алиев~Р.~А.

\paper Теорема единственности одной обратной задачи для квазилинейного уравнения эллиптического типа

\serial Матем. моделирование и краев. задачи

\yr 2010

\pages 13--15

\mathnet{http://mi.mathnet.ru/mmkz1626}

\inbook Труды седьмой Всероссийской научной конференции с~международным участием. \rm Часть~3

\bookinfo Дифференциальные уравнения и краевые задачи

\yr 2010

\publ СамГТУ

\publaddr Самара





\RBibitem{ali:bib:24} 

\by Ладыженская~О.~Г., Уральцева~Н.~Н.

\book Линейные и квазилинейные уравнения эллиптического типа

\publaddr М.

\publ Наука 

\yr 1973 

\totalpages 576

\zmath{https://zbmath.org/?q=an:0269.35029}

%Ladyzhenskaya, O.A.; Ural’tseva, N.N.

%Linear and quasilinear equations of elliptic type. 2nd revised ed. (Линейные и квазилинейные уравнения эллиптического типа.) (Russian) Zbl 0269.35029

%Moskva: Izdat. ”Nauka”. 576 p. R. 2.50 (1973).





\Bibitem{ali:bib:25} 

\zmath{https://zbmath.org/?q=an:0198.14101}

\by Miranda~C.

\book Partial differential equations of elliptic type

\publaddr Berlin, Heidelberg, New York

\publ Springer Verlag 

\totalpages xii+370 

\yr 1970



\end{thebibliography}



\RuDate



\makeentitle



%\thanks{{\bf ORCID}

%

%{\bf Ramiz Aliyev:} \url{http://orcid.org/xxxx-xxxx-xxxx-xxxx}

%}







\EngRef



\begin{thebibliography}{99}







\Bibitem{ali:bib:1-eng} 

\lang In Russian

\by Lavrent'ev~M.~M., Romanov~V.~G., Shishatskii~S.~P.

\book Nekorrektnye zadachi matematicheskoi fiziki i analiza \rm [Ill-posed problems of mathematical physics and analysis]

\publaddr Moscow

\publ Nauka 

\yr 1980

\totalpages 288





\Bibitem{ali:bib:2-eng} 

\lang In Russian

\by Kabanikhin~S.~I.

\book Obratnye i nekorrektnye zadachi \rm [Inverse and ill-posed problems]

\publaddr Moscow

\publ Nauka 

\yr 2009

\totalpages 458









\Bibitem{ali:bib:4-eng} 

\lang In Russian

\by Iskenderov~A.~D.

\paper Some inverse problems on the determination of right-hand parts of differential equations

\jour Izv. Akad. Nauk Az. SSR

\yr 1976

\issue 2

\pages 58--63 

\yr 1976





\Bibitem{ali:bib:5-eng} 

\lang In Russian

\by Iskenderov~A.~D.

\paper Inverse problem about the determination of the coefficients of a quasilinear elliptic equation

\jour  Izv. Akad. Nauk Az. SSR

\yr 1978

\issue 2

\pages 80--85 

\yr 1978





\Bibitem{ali:bib:6-eng} 

\zmath{https://zbmath.org/?q=an:0428.35083}

\by Iskenderov~A.~D.

\paper The inverse problem of the determination of the coefficients in an elliptic equation

\jour Differ. Equ.

\vol 15

\pages 605--612 

\yr 1979





\Bibitem{ali:bib:7-eng} 

\lang In Russian

\by Temirbulatov~S.~I.

\book Obratnye zadachi dlia ellipticheskikh uravnenii \rm [Inverse problems for elliptic equations]

\publaddr Alma-Ata

\publ Kazakh. Univ.

\yr 1975 

\totalpages 72







\Bibitem{ali:bib:8-eng} 

\lang In Russian

\by Klibanov~M.~V.

\paper A two-dimensional inverse problem for an elliptic equation

\jour Differ. Uravn. 

\vol 19

\issue 6

\pages 1072--1074 

\yr 1983







\Bibitem{ali:bib:9-eng} 

\by Klibanov~M.~V.

\zmath{https://zbmath.org/?q=an:0573.35083}

\paper Inverse problems in the ``large'' and Carleman bounds

\jour Differ. Equ.

\vol 20

\issue 6

\pages 755--760 

\yr 1984





\Bibitem{ali:bib:10-eng} 

\zmath{https://zbmath.org/?q=an:0567.35080}

\by Klibanov~M.~V.

\paper Uniqueness in the large of solutions of inverse problems for a class of differential equations

\jour Differ. Equ.

\vol 20

\issue 11

\pages 1390--1395 

\yr 1984





\RBibitem{ali:bib:11} 

\lang In Russian

\by Khaidarov~A. 

\yr 1984

\vol 277

\issue 6 

\pages 1335--1337

\paper One class of inverse problems for elliptic equations 

\jour Dokl. Akad. Nauk SSSR









\Bibitem{ali:bib:12-eng} 

\lang In Russian

\inbook Nekotorye zadachi matema\-tiches\-koi fiziki i analiza \rm [Some Problems of Mathematical Physics and Analysis]

\by Khaidarov~A. 

\paper On an inverse problem for elliptic equations

\publ Nauka, Siberian Branch

\publaddr Novosibirsk 

\yr 1984

\pages 245–249





\Bibitem{ali:bib:13-eng} 

\by Vabishchevich~P.~N.

\paper An inverse problem for reconstructing the right-hand side of an elliptic equation and its numerical solution

\zmath{https://zbmath.org/?q=an:0595.65130}

\jour Differ. Equ.

\vol 21

\pages 201--207 

\yr 1985



\Bibitem{ali:bib:14-eng} 

\zmath{https://zbmath.org/?q=an:0695.35214}

\by Vabishchevich~P.~N.

\paper Uniqueness of solutions of some inverse problems for elliptic equations

\jour Differ. Equ.

\vol 24

\issue 12

\pages 1443--1446 

\yr 1988





\Bibitem{ali:bib:15-eng} 

\by Solov'ev~V.~V. 

\zmath{https://zbmath.org/?q=an:1160.35549}

\paper Inverse problems for elliptic equations on the plane.~I.

\jour Differ. Equ. 

\vol 42

\issue 8

\pages 1170--1179 

\yr 2006

\crossref{10.1134/s0012266106080118}

\scopus{2-s2.0-33750032371}



\Bibitem{ali:bib:16-eng} 

\by Solov'ev~V.~V. 

\paper Source and coefficient inverse problems for an elliptic equation in a rectangle

\jour Comput. Math. Math. Phys.

\yr 2007

\vol 47

\issue 8

\pages 1310--1322

\mathscinet{http://www.ams.org/mathscinet-getitem?mr=2378182}

\crossref{10.1134/S096554250708009X}

\scopus{2-s2.0-34548411765}





\Bibitem{ali:bib:17-eng} 

\paper A Global Uniqueness Theorem for an Inverse Boundary Value Problem

\jour Annals of Mathematics 

\vol 125 

\issue 1 

\pages 153--169

\by Sylvester~J., Uhlmann~G.

\crossref{10.2307/1971291}



\Bibitem{ali:bib:18-eng} 

\paper Inverse coefficient problems for nonlinear elliptic equations

\jour ANZIAM Journal 

\vol 49 

\yr 2008

\issue 02 

\pages  271--279

\by Yang~R., Ou~Y.

\crossref{10.1017/s1446181100012839}

\scopus{2-s2.0-77954661742}





\Bibitem{ali:bib:19-eng} 

\lang In Russian

\by Vakhitov~I.~S.

\paper Inverse problem of identification of the diffusion coefficient in diffision-reaction equation

\jour Dal'nevost. matem. zhurn.

\yr 2010

\vol 10

\issue 2

\pages 93--105







\Bibitem{ali:bib:20-eng} 

\by Kozlov~V.~A., Maz'ya~V.~G., Fomin~A.~V.

\paper An iterative method for solving the Cauchy problem for elliptic equations

\jour U.S.S.R. Comput. Math. Math. Phys.

\yr 1991

\vol 31

\issue 1

\pages 45--52

\mathscinet{http://www.ams.org/mathscinet-getitem?mr=1099360}

\isi{http://gateway.isiknowledge.com/gateway/Gateway.cgi?GWVersion=2&SrcApp=PARTNER_APP&SrcAuth=LinksAMR&DestLinkType=FullRecord&DestApp=ALL_WOS&KeyUT=A1991HU21600006}





\Bibitem{ali:bib:21-eng} 

\paper An iterative procedure for solving a Cauchy problem for second order elliptic equations

\yr 2004 

\jour Mathematische Nachrichten 

\vol 272 

\issue 1 

\pages 46--54

\by Johansson~T.

\crossref{10.1002/mana.200310188} 

\scopus{2-s2.0-4644351704}







\Bibitem{ali:bib:22-eng} 

\lang In Russian

\by Aliyev~R.~A.

\paper The uniqueness theorem of an inverse problem for a quasilinear elliptic equation

\jour Izvestiia vuzov. Severo-Kavkazskii region. Estestvennye nauki

\yr 2012

\issue 5 

\pages 5--10





\Bibitem{ali:bib:23-eng}

\lang In Russian

\by Aliyev~R.~A.

\paper The uniqueness theorem of an inverse problem for a quasilinear elliptic equation

\serial Matem. modelirovanie i kraev. zadachi

\yr 2010

\pages 13--15

\yr 2010

\publ Samara State Technical Univ.

\publaddr Samara





\Bibitem{ali:bib:24-eng} 

\lang In Russian

\yr 1973 

\totalpages 576

\zmath{https://zbmath.org/?q=an:0269.35029}

\by Ladyzhenskaya~O.~A., Ural'tseva~N.~N.

\book Lineinye i kvazilineinye uravneniia ellipticheskogo tipa \rm [Linear and quasilinear equations of elliptic type]

\publaddr Moscow

\publ Nauka







\Bibitem{ali:bib:25-eng} 

\zmath{https://zbmath.org/?q=an:0198.14101}

\by Miranda~C.

\book Partial differential equations of elliptic type

\publaddr Berlin, Heidelberg, New York

\publ Springer Verlag 

\totalpages xii+370 

\yr 1970



\end{thebibliography}



\EngDate





\renewcommand{\baselinestretch}{0.9}





\newpage \Russian



%\end{document}

%

%

